\section{Background}

The capacity of a lithium-ion battery --- the total charge deliverable in a full discharge --- fades 
with repeated cycling. With time and use, storage capacity diminishes and internal resistance increases, driven by 
degradation mechanisms that can occur simultaneously or trigger one another \cite{edge2021}. A battery is said to 
reach end-of-life (EoL) when its capacity drops below 80\% of its initial value, following the criterion used by 
Zhou et al.\ \cite{ZHOU20236117}.

The rate at which a battery charges or discharges is characterized by the C-rate, defined as the ratio of applied 
current to nominal capacity \cite{ossila_crate}. At 1C, a full charge or discharge completes in one hour; at 2C, 
in thirty minutes. Higher C-rates induce greater mechanical stress during cycling and promote lithium plating during 
charge \cite{edge2021}, both of which accelerate capacity fade. In the dataset used in this study, all cells are 
discharged at a fixed 4C rate while the charging protocol varies --- making the average charging C-rate the primary 
source of variability in operating conditions across cells.